\section{Riesz-weighted prime sums and zero detection on windows}\label{sec:landau}

\subsection{Weighted prime sums}

\begin{definition}\label{def:landau:G}
For $y\in\R$ put
\[
G(y)=\sum_{1\le n\le y}\Lambda(n)\max\Big(\frac yn-1,0\Big)-y\,(\log y-1-\gamma),
\]
and for a Dirichlet character $\chi\ne\one$
\[
g_\chi(y)=\sum_{1\le n\le y}\Lambda(n)\chi(n)\max\Big(\frac yn-1,0\Big)+y\,\ell(\chi),
\]
with $g_{\one_1}=G$. \formal{Formalism.HC.An.G} \formal{Formalism.HC.An.gChi}
\end{definition}

Using $\int_1^\infty\max(y/n-1,0)\,y^{-s-1}\,dy=n^{-s}/(s(s-1))$ one finds, for $\Re s>1$,
\begin{equation}\label{eq:landau:mellin}
\int_1^\infty g_\chi(y)\,y^{-s-1}\,dy=-\frac{L'/L(s,\chi)-s\,\ell(\chi)}{s(s-1)}\qquad(\chi\ne\one),
\end{equation}
and similarly for $G$ with $\zeta'/\zeta(s)+\frac1{s-1}+1-(1+\gamma)s$ in the numerator. The right-hand
side is regular at $s=1$ and has a simple pole with residue $-m/(\rho(\rho-1))$ at each zero $\rho\neq0$
of multiplicity $m$. Heuristically $g_\chi(y)=-\sum_\rho y^\rho/(\rho(\rho-1))+\cdots$, so GRH is
equivalent to $g_\chi(y)\ll y^{1/2+\eps}$. No explicit formula with a sum over zeros is used below.

\begin{definition}\label{def:landau:Gm}
Let $\psi\in\mathcal P$, let $o$ be the order of $\chi_\psi$, and let $\psi^k=(\chi_\psi^k)^{\rm prim}$
for $0\le k<o$, so that $\psi^0=\one$. Put
\[
G_\psi(y)=G(y)+\frac1{3o+1}\sum_{k=1}^{o-1}\Re g_{\psi^k}(y).
\]
\formal{Formalism.HC.Level.Gm} (cyclic level \formal{Formalism.HC.Level.cyc})
\end{definition}

The weight $(3o+1)^{-1}$ is the common value of the non-trivial Fourier coefficients of the coset
distribution used in \cref{sec:hosts} (\cref{rem:hosts:fourier}). It matters only that the weights are
positive.

\subsection{A Landau-type theorem on windows}

A \emph{character family} $F$ is a finite list of pairs $(\chi_i,a_i)$, where $\chi_i$ is a character
modulo $N_i$ with $N_i=1$ or $\chi_i\ne\one$, and $a_i>0$. Put
$\mathrm{Gm}_F(y)=\sum_ia_i\Re g_{\chi_i}(y)$. Thus $G_\psi=\mathrm{Gm}_F$ for the family $F$ of the
$\psi^k$, with weights $1$ and $(3o+1)^{-1}$.

\begin{theorem}[Landau's theorem on windows]\label{thm:landau-window}
Let $F$ be a character family, $0<\theta<1$, $\delta>0$ and $C\in\R$. Suppose that for every $Y_0$
there is $Y\ge Y_0$ such that
\[
\mathrm{Gm}_F(y)\le C\,y^\theta\qquad\text{for all }y\in[Y,Y^{1+\delta}].
\]
Then $L(s,\chi_i)\ne0$ for every $i$ and every $s$ with $\Re s>\theta$.
\formal{Deduction.lemmaL} \formal{Deduction.landau_window}
\end{theorem}

The classical theorem of Landau \cite{Landau1905} (see also \cite{Ingham1932}) gives the same
conclusion if $\mathrm{Gm}_F(y)\le Cy^\theta$ for \emph{all} large $y$. The point of
\cref{thm:landau-window} is that the one-sided bound is only required on a sequence of windows of
fixed multiplicative length.

\begin{proof}
\emph{Step 1: Laplace transform and poles.} Put $g(u)=\mathrm{Gm}_F(e^u)$. Then $g$ is measurable and
$|g(u)|\le C_0e^{2u}$ for $u\ge0$. For $\Re s>2$ let $\Phi(s)=\int_0^\infty g(u)e^{-su}\,du$.

For each $i$ let $f_i(s)=L(s,\chi_i)$ if $\chi_i\ne\one$ and $f_i(s)=(s-1)\zeta(s)$ if $N_i=1$. Then
$f_i$ is entire, and $f_i(1)\ne0$. By \eqref{eq:landau:mellin} and the identity
$\overline{g_\chi(y)}=g_{\bar\chi}(y)$,
\[
\Phi(s)=\sum_i\frac{a_i}2\big(\Phi_{\chi_i}(s)+\Phi_{\bar\chi_i}(s)\big),\qquad
\Phi_\chi(s)=-\frac{f'/f(s)-s\,f'/f(1)+k_\chi(1-s)}{s(s-1)}
\]
for suitable constants $k_\chi$. Hence $\Phi$ continues meromorphically to $\Re s>0$. Its poles form a
discrete set $Z$, namely the zeros $\rho$ of the $L(s,\chi_i)$ with $\Re\rho>0$. Each pole is simple, with
residue
\[
-\frac1{\rho(\rho-1)}\sum_i\frac{a_i}2\big(m_i(\rho)+m_i(\bar\rho)\big)\neq0,
\]
where $m_i$ is the multiplicity. At a real pole $\beta\in(0,1)$ the residue is \emph{positive}, because
$a_i>0$ and $-1/(\beta(\beta-1))>0$. There are no zeros with $\Re s\ge1$, so $\Re\rho<1$ on $Z$.

\emph{Step 2: Taylor coefficients.} For $\Re s>2$ put $a_k(s)=(-1)^k\Phi^{(k)}(s)/k!$. Differentiating
under the integral,
\[
a_k(s)=\int_0^\infty g(u)K_{k,s}(u)\,du,\qquad K_{k,s}(u)=\frac{u^ke^{-su}}{k!}.
\]
Since $|K_{k,\sigma+i\gamma}|=K_{k,\sigma}\ge0$, a one-sided bound $g\le R$ (with $R\ge0$) on an interval
$[U,V]$ gives, for every $\gamma\in\R$ and $\sigma>2>\theta$,
\begin{equation}\label{eq:landau:onesided}
|a_k(\sigma+i\gamma)|+\Re a_k(\sigma)\le2\int g^+K_{k,\sigma}
\le\frac{2}{(\sigma-\theta)^k}\int_U^VR(u)e^{-\theta u}du+2C_0\!\int_{(0,\infty)\smallsetminus[U,V]}\!\frac{u^ke^{(2-\sigma)u}}{k!}du.
\end{equation}
This uses $u^ke^{-(\sigma-\theta)u}/k!\le(\sigma-\theta)^{-k}$, and $g^+\le|g|$ outside $[U,V]$.

\emph{Step 3: the key inequality.} A Chernoff-type estimate shows that for every $\delta>0$ there are
$\sigma>2$ and $C_1$ such that for every $U\ge1$ there is an integer $k\ge U$ with
$\int_{(0,\infty)\smallsetminus[U,(1+\delta)U]}u^ke^{(2-\sigma)u}/k!\,du\le C_1\sigma^{-k}$. The kernel
concentrates near $u=k/(\sigma-2)$.

By hypothesis, for every $\eta>0$ and $U_0$ there are $U\ge U_0$ and $R(u)=|C|e^{\theta u}$ with $g\le R$
on $[U,(1+\delta)U]$; then $\int_U^{(1+\delta)U}Re^{-\theta u}du=|C|\delta U\le e^{\eta U}$ for large $U$.
Inserting this into \eqref{eq:landau:onesided} gives the following: for every $\eta>0$ and every $U_0$
there are $U\ge U_0$ and an integer $k\ge U$ such that, with $C_2=2+2C_0C_1$,
\begin{equation}\label{eq:landau:key}
|a_k(\sigma+i\gamma)|+\Re a_k(\sigma)\le C_2\,e^{\eta U}(\sigma-\theta)^{-k}\qquad(\gamma\in\R).
\end{equation}

\emph{Step 4: the contradiction.} Suppose some $\rho_0\in Z$ has $\Re\rho_0>\theta$. We use three
standard facts.
\begin{itemize}
\item \emph{Pole asymptotics.} If $\Phi$ is holomorphic on $B(z,r)\smallsetminus\{\rho\}$, where $\rho$
is a simple pole with residue $c$, and $r'<r$, then $|a_k(z)-c(z-\rho)^{-k-1}|\le M r'^{-k}$.
\item \emph{Cauchy's estimate.} If $\Phi$ is holomorphic on a disc of radius $r>r'$ about $z$, then
$|a_k(z)|\le Mr'^{-k}$.
\item \emph{Growth comparison.} If $0<d<d'\le E$ and $c>0$, it is impossible that for every $\eta>0$
there are arbitrarily large $U$ and $k\ge U$ with $c\,d^{-k}\le M{d'}^{-k}+C_2e^{\eta U}E^{-k}$. To see
this, take $\eta=\frac12\log(d'/d)$.
\end{itemize}
The set $Z\cap\overline B(\sigma,\sigma-\theta)$ is finite and non-empty, since it contains $\rho_0$.
Hence one of the following two cases occurs.

(A) The point of $Z$ nearest to $\sigma$ is not real, or there are several nearest points. Moving
the centre to a suitable height $\sigma+i\gamma$ produces $r\le\sigma-\theta$ and $\rho\in Z$ such
that $\rho$ is the unique point of $Z$ in $B(\sigma+i\gamma,r)$. Moreover $\Phi$ is holomorphic on a
disc about $\sigma$ of radius larger than $d'=(|\sigma+i\gamma-\rho|+r)/2$.

By pole asymptotics, $|a_k(\sigma+i\gamma)|\ge|c|d^{-k-1}-Md'^{-k}$ with $d=|\sigma+i\gamma-\rho|$.
By Cauchy's estimate, $\Re a_k(\sigma)\ge-M'd'^{-k}$. Inserting both into \eqref{eq:landau:key}
contradicts the growth comparison with $E=\sigma-\theta$.

(B) The nearest point is real and isolated, say $\beta\in(\theta,1)$, with residue $c>0$. Pole
asymptotics give $\Re a_k(\sigma)\ge c(\sigma-\beta)^{-k-1}-Md'^{-k}$, where
$\sigma-\beta<d'\le\sigma-\theta$. Together with $|a_k(\sigma)|\ge0$ and \eqref{eq:landau:key} at
$\gamma=0$, this again contradicts the growth comparison. It is here that the positivity of the
residues at real poles replaces the positivity of the function in Landau's argument.

Hence every zero $\rho$ of every $L(s,\chi_i)$ with $\Re\rho>0$ satisfies $\Re\rho\le\theta$.
\end{proof}

\begin{corollary}\label{cor:landau:formB}
Let $\psi\in\mathcal P$, $\lambda>1$ and $\theta\in(0,1)$. If $\{x:G_\psi(x)\le x^\theta\}$ has
windows of exponent $\lambda$, then $\zeta(s)$ and $L(s,\psi^k)$, $0<k<o$, have no zeros with
$\Re s>\theta$. In particular, $\psi$ is zero-free above $\theta$.
\formal{UpperTheory.Comparison₂.formB}
\end{corollary}

\begin{proof}
Apply \cref{thm:landau-window} to the family of \cref{def:landau:Gm}, with $\delta=\lambda-1$ and
$C=1$. Note that $\psi=\psi^1$.
\end{proof}
